\chapter{化学方程式是反应的账本}

\begin{kaichang}
过生日的时候，点上一支蜡烛。蜡烛越烧越短，最后只剩一小摊蜡油。蜡去哪了？很多人会说：“烧掉了呀。”可是“烧掉”到底是什么意思——是凭空消失了吗？

再看另一件事。把一枚铁钉放在潮湿的窗台上，几周后它长满了红棕色的铁锈。如果把生锈前后的铁钉都拿去称一称，你会发现一件怪事：生锈之后的钉子，比原来\textbf{变重了}。

一边在“消失”，一边在“变重”。物质到底守不守规矩？这一章，我们要给化学反应立一本账本——这本账本的名字，你已经很熟悉了：化学方程式。
\end{kaichang}

\section{蜡去哪了：先称一称}

先别急着回答“蜡去哪了”，我们像侦探一样，只做能测量的事。

找一支新蜡烛，称出质量，点燃，让它烧掉一半，再称。结果毫无悬念：变轻了。如果把一块干冰（固态二氧化碳）放在敞口盘子里，它会越变越小，最后“消失”；把浓盐酸敞口放一会儿，瓶子的总质量也会变轻。看起来，物质真的会跑掉。

但请等一下。上面这些实验都有一个共同点：容器是\textbf{敞开的}。气体可以进来，也可以出去。如果称的不是“反应前后那几样东西”，而是“房间里所有东西的总账”，结果会怎样？

生活里这样的“假失踪”到处都是。汽油烧完，油箱空了，可一辆车每烧掉 \ud{1}{L} 汽油，会向空气放出 \ud{2}{kg} 多的二氧化碳——车变轻了，空气变“重”了。食物变质发臭，是有机分子和氧气、微生物里的物质重新组合，一部分产物散进了空气。切开的苹果放久了变褐，也是果肉里的物质和氧气悄悄反应。反过来，“凭空变多”的账也一样：小树苗长成大树，增加的质量绝大部分不是从土壤里“吃”来的，而是从空气里的二氧化碳和水变来的。称不清，往往是因为边界没画对。

\begin{shiyan}
\textbf{密闭瓶里的冒泡反应}

材料：小苏打（碳酸氢钠，食用级）、白醋、一个带盖的塑料瓶（或锥形瓶）、一只气球、一台厨房电子秤（精确到 \ud{0.1}{g} 更好）。

步骤：
\begin{enumerate}[itemsep=1pt]
\item 往瓶里倒入约 \ud{50}{mL} 白醋；往气球里装两小勺小苏打。
\item 把气球口套在瓶口上，注意先别让小苏打掉进去。把整个装置（瓶+气球+里面的全部东西）放到电子秤上，记下总质量。
\item 提起气球，让小苏打全部落入醋中。瓶里立刻冒泡，气球慢慢鼓起来。
\item 等反应平息后，再称一次整个装置的总质量，与第 2 步比较。
\end{enumerate}

观察点：气球鼓起来，说明有气体生成；但整个装置的总质量几乎不变（差别通常在秤的误差范围内）。

安全提示：这个实验用的是厨房材料，可以在家做，但请在通风处操作，做完洗手，瓶内液体不要入口；气球别套得太紧，防止被胀破。更剧烈的密闭反应（如燃烧的密闭容器实验）请看老师演示或视频，不要自己尝试。
\end{shiyan}

气泡明明产生了，气体明明胀大了气球，总质量却纹丝不动。这说明什么？说明那个反应里，没有任何一份物质凭空出现或凭空消失——小苏打和醋里的“料”，只是换了一种存在形式，一部分变成了气球里的二氧化碳气体。如果你把容器敞开，气体跑掉，秤就会“撒谎”说质量变少了；把门窗关严，账就平了。

现在回头看蜡烛。蜡烧掉后生成的主要是二氧化碳和水蒸气，都是看不见的气体，悄悄散进了空气。如果你把蜡烛和燃烧需要的空气一起封在一个大玻璃罩里称，总质量依然守恒。而铁钉生锈变重，是因为空气里的氧气“入伙”了——多出来的质量，正是参加反应的氧原子的质量（第 15 章讲过铁锈的故事）。

\section{原子不会消失，只会换搭档}

现象层的结论很清楚：\textbf{反应前后，总质量不变}。可为什么不变？把镜头推进到粒子层，答案简单得漂亮。

回忆第 6 章：道尔顿的原子论是被定比定律、倍比定律这些称量证据“逼”出来的。化学反应发生时，原子本身既不诞生，也不毁灭，也不会从一种元素变成另一种元素（那需要动原子核，是第 8 章核反应的地盘，普通化学反应的能量根本够不着）。原子做的事只有一件：\textbf{重新组合}。

拿氢气在氧气中燃烧来说。反应前，是氢分子 \chem{H_2} 和氧分子 \chem{O_2}——各自两个原子抱在一起。反应时，旧搭档拆开（断开旧的化学键需要吸收能量），原子们重新配对，组成新的水分子 \chem{H_2O}（形成新键会放出能量；这笔能量账怎么算，是下一章的主题）。整场变化里，氢原子还是氢原子，氧原子还是氧原子，一个不多，一个不少。

所以质量守恒根本不是什么神秘的定律，它只是“原子守恒”在秤上的投影：每种原子的数目不变，每个原子的质量不变，总质量自然不变。反过来说也一样硬气：如果你在哪个实验里发现总质量真的变了，那只有两种可能——有东西漏进或漏出了你的“账本边界”，或者你遇上了核反应。查账先查边界，这是化学家的第一反应。

把“重新组合”这四个字再放大一点看。你可以把一场化学反理想象成一场大型舞会换舞伴：音乐响起，原来的成对舞者松开手，在人群里转一圈，又按新的搭配牵起手。整场舞会，人头数从头到尾没变，变的只是“谁和谁在一起”。燃烧、生锈、冒泡、变色，千奇百怪的现象底下，都是同一件事：原子换搭档。而舞会为什么这样换而不是那样换——谁和谁的组合能量更低、换起来有多快——那是后面几章的规则。这一章先立账：人头必须点对。

\section{配平：让每类原子两边相等}

既然每种原子的数目在反应前后必须相等，我们就可以把这个约束写成规则。先把氢气燃烧用化学式摆出来：

\[\chem{H_2} + \chem{O_2} \rightarrow \chem{H_2O}\]

数一数：左边 2 个氢原子、2 个氧原子；右边 2 个氢原子、却只有 1 个氧原子。账不平。怎么办？

关键纪律只有一条：\textbf{只许改化学式前面的系数，绝不许动化学式里的下标}。下标是分子的“配方”，\chem{H_2O} 改成 \chem{H_2O_2}，那已经不是水，而是双氧水了——等于为了平账，把账本上的货物偷换成了别的东西。

于是正确的做法是调整“取几份”：取 2 份氢分子，生成 2 份水分子，氧原子正好用掉 2 个：

\[\chem{2H_2} + \chem{O_2} \rightarrow \chem{2H_2O}\]

再数一遍：左边 4 个氢、2 个氧；右边 4 个氢、2 个氧。平了。这个过程叫\textbf{配平}。它不是数学游戏，不是“把数字凑得好看”，而是把“每种原子总数不变”这条粒子层的事实，翻译成符号层的等式。

再来一个稍复杂的。甲烷（天然气的主要成分）燃烧：

\[\chem{CH_4} + \chem{O_2} \rightarrow \chem{CO_2} + \chem{H_2O}\]

碳原子左右各 1 个，已经平了；氢左边 4 个，右边只有 2 个，给水配上系数 2；这样右边氧原子共 4 个，给氧气配上系数 2：

\[\chem{CH_4} + \chem{2O_2} \rightarrow \chem{CO_2} + \chem{2H_2O}\]

配平的手感来自练习，但心法永远不变：一边数原子，一边调份数，最后所有系数约成最简整数比。配平之后，系数告诉你的是粒子的\textbf{个数比}——每 1 个甲烷分子消耗 2 个氧分子。而“个数比”通过第 5 章的摩尔桥梁，立刻就能换算成“质量比”，这正是化学家配料的依据（本章后面会算给你看）。

再看一个带条件的例子。双氧水（过氧化氢溶液）放着会慢慢分解，滴一点二氧化锰粉末进去，立刻剧烈冒泡：

\[\chem{2H_2O_2(aq)} \xrightarrow{\text{催化剂}} \chem{2H_2O(l)} + \chem{O_2(g)}\]

配平过程你自己可以走一遍：左边 4 个氢、4 个氧，右边先配 2 个水用掉 4 个氢和 2 个氧，剩下 2 个氧恰好凑成 1 个氧分子。箭头上方的“催化剂”三个字不参与数数——二氧化锰在反应前后本身没有被消耗，它只是让这条路走得更快（第 30 章细讲它的工作方式），所以账本上不需要给它留位置。这个反应也是实验室制氧气的常用办法，冒出的气泡能让带火星的木条复燃。

\begin{wuqu}
“配平就是为了把数字凑整齐，实在配不平就改一下化学式。”——这是最危险的一条捷径。有同学配氢气燃烧时，左边氧原子多一个，就把水写成 \chem{H_2O_2}。账是“平”了，物质却被换掉了：双氧水和水是两种性质完全不同的物质（一个能漂白消毒，一个是你每天喝的东西）。方程式里的每个化学式都对应一种真实存在的物质，化学式不是可以随便涂改的草稿。配不平，只能说明你设的份数不对，或者反应物、生成物写错了——账平不了时，该查的是货，不是改账本。
\end{wuqu}

\section{一行方程式里藏着的全部信息}

现在把符号层补完整。一条规范的化学方程式，提供的信息远不止“谁变成谁”：

\begin{itemize}[itemsep=2pt]
\item \textbf{反应物}写在箭头左边，\textbf{生成物}写在右边，箭头读作“反应生成”。
\item \textbf{系数}给出各物质的粒子个数比（也就是物质的量之比），由原子守恒配平而来。
\item \textbf{状态符号}跟在化学式后面的小括号里：固态 (s)、液态 (l)、气态 (g)、水溶液 (aq)（aqueous，“含水的”——第 25 章讲过，溶解的粒子在水里的处境和纯物质很不一样）。同一件东西，状态不同，反应的表现可以完全不同：气态氢气遇火星就爆燃，氢气溶在水里却安静得很。
\item \textbf{反应条件}写在箭头上方或下方。最常见的是加热（常写成 $\Delta$）；如果反应靠催化剂帮忙，就把催化剂写在箭头上方——注意，催化剂只出现在“过路”的位置，不进入总账：它提供了另一条更好走的路，自己却不在最终的物料清单里（第 30 章会专门讲它怎样另辟蹊径）。
\item \textbf{可逆箭头} $\rightleftharpoons$：有些反应不是单行道，生成物还能变回去，两边同时进行。这时方程式的账要朝两个方向一起记——哪边占上风、能不能被条件扳动，是第 31 章化学平衡的主题，这里先认得这个符号就好。
\end{itemize}

举个把信息写全的例子。煅烧石灰石（碳酸钙）制生石灰，工业上在大约 $\ud{900}{^{\circ}C}$ 的高温炉里进行：

\[\chem{CaCO_3(s)} \xrightarrow{\text{高温}} \chem{CaO(s)} + \chem{CO_2(g)}\]

读懂这一行，你就知道了：固体碳酸钙在高温下分解，生成固体氧化钙和二氧化碳气体，三者的粒子数之比是 $1:1:1$。石灰窑里气体跑掉，剩下的固体质量变轻——账依然平，只是有一部分“货”从烟囱离开了账本的边界。

\begin{zhengju}
“总质量不变”这条规则不是哪位圣人宣布的，它是被一台天平一点点逼出来的。18 世纪的化学被“燃素说”统治着：人们相信可燃物里含有一种叫“燃素”的东西，燃烧就是燃素跑掉。这套说法能解释“木头烧完变轻”，却在“金属煅烧后变重”面前狼狈不堪——支持者只好说燃素的质量是负的。一个理论开始靠打补丁续命，通常离翻船不远了。

1770 年代，法国化学家拉瓦锡做了一系列以\textbf{称量}为核心的实验。最著名的一个：他把锡放在密闭的玻璃容器里加热。锡表面生成了灰白色的“煅灰”，但整个密闭容器的总质量不变。打开容器时，空气“嘶”地冲了进去——说明容器里有一部分空气被用掉了。再称：锡变成煅灰所增加的质量，恰好等于容器里少掉的那部分气体的质量。一来一回，分毫不差。

替代解释当场被拆穿：如果燃烧真是“燃素逃逸”，密闭容器的总质量理应变轻，可它没有；而“金属从空气中拿走了某种东西”这个解释，同时说清了变重的来源和冲进去的空气。拉瓦锡后来用汞做了更精细的版本，还精确测出了空气里被消耗的那部分约占五分之一——也就是后来被命名为氧气的成分（第 14 章）。

这个故事值得记住的有两点。第一，拉瓦锡的“秘密武器”不是天才直觉，而是把天平用到极致：他把“密闭”当作账本的边界，把每一次质量变化都追到一个具体的去向。第二，质量守恒最初也只是大量实验的归纳，没人见过原子——直到一百多年后，原子论确立（第 6 章），这条经验定律才获得了粒子层的解释：“物质不灭”其实是“原子不换”。好的科学规律就是这样：先有用，后有解，解释来了，规律反而站得更稳。
\end{zhengju}

\section{按账本配料：化学计量}

配平好的方程式不只是“谁变成谁”的叙述，它还是一张\textbf{配料表}。系数之比 $=$ 粒子数之比 $=$ 物质的量之比。搭上第 5 章的摩尔桥梁，质量也能算。

\begin{qiaoliang}
氢气燃烧：$\chem{2H_2} + \chem{O_2} \rightarrow \chem{2H_2O}$。系数说：每 2 份 \chem{H_2} 配 1 份 \chem{O_2}，生成 2 份 \chem{H_2O}。换成摩尔：\ud{2}{mol} 氢气配 \ud{1}{mol} 氧气，生成 \ud{2}{mol} 水。再换成质量（用第 5 章的相对分子质量）：氢气 $2\times 2=\ud{4}{g}$，氧气 $\ud{32}{g}$，生成水 $2\times 18=\ud{36}{g}$。验算：$4+32=36$，总质量分毫不差——账本的借贷两边永远相等。火箭工程师给发动机配燃料，算的就是这样的账：液氢和液氧按质量 $1:8$ 上箭，不是拍脑袋，是从这一行方程式里读出来的。
\end{qiaoliang}

现实中更常见的问题是：配料不可能总是刚刚好。谁先用完，反应就停在哪儿——先用完的那个叫\textbf{限制反应物}。

先用生活里的例子找感觉。你做三明治：每份需要 2 片面包夹 1 片奶酪。现有 10 片面包、4 片奶酪，能做几份？奶酪只够 4 份，用掉 8 片面包，还剩 2 片面包干瞪眼。奶酪是限制反应物，它决定了产量；面包是过量的。

化学里一模一样。把 \ud{4}{g} 氢气（\ud{2}{mol}）和 \ud{16}{g} 氧气（\ud{0.5}{mol}）放在一起点燃：按 $2:1$ 的配比，\ud{2}{mol} 氢气需要 \ud{1}{mol} 氧气，可氧气只有 \ud{0.5}{mol}，只够烧掉 \ud{1}{mol} 氢气。结果：氧气用尽，生成 \ud{1}{mol} 水（\ud{18}{g}），还剩 \ud{1}{mol} 氢气原封不动。不管你把它们关在一起多久，剩下的氢气也不会自己变成水——巧妇难为无米之炊。

现在回过头，给开头那个气球实验算算账。小苏打和醋的反应是

\[\chem{NaHCO_3} + \chem{CH_3COOH} \rightarrow \chem{CH_3COONa} + \chem{H_2O} + \chem{CO_2}\]

系数全是 1：每 1 份小苏打生成 1 份二氧化碳。小苏打的摩尔质量约 \ud{84}{g/mol}，两小勺大约 \ud{5}{g}，也就是约 \ud{0.06}{mol}，所以理论上生成约 \ud{0.06}{mol} 二氧化碳。常温常压下 \ud{1}{mol} 气体大约占 \ud{24}{L}，$0.06\times 24\approx\ud{1.4}{L}$——你的气球应该鼓成一个直径十几厘米的球。下次做的时候，不妨先算后做，看看气球的大小和账本上预计的差多少。这种“先估数量级再动手”的习惯，是化学家避免出错的护身符：如果气球鼓成了篮球那么大，或者纹丝不动，你立刻就知道哪里出了问题。

医院里也在算同样的账。医生开药常按体重计算剂量——每公斤体重多少毫克——因为药物分子要和你体内的物质按粒子的比例打交道；透析液、输液里各种离子的浓度，更是照着化学计量的精确配比调出来的，错一个数量级就可能出事故。化学计量不是考试技巧，是所有“按比例办事”的通用语言。

如果反应物的量都按方程式算准了，理论上能得到的最大产量，叫\textbf{理论产率}。可你走进任何一家化工厂或实验室问一问：真的有人拿到过 100\% 吗？几乎没有。实际产量除以理论产量，叫\textbf{产率}。它为什么会“打折”？账本之外的原因有一串：

\begin{itemize}[itemsep=2pt]
\item \textbf{副反应}：原料可能同时走了另一条路，生成你不想要的东西（高温下碳燃烧既可能生成 \chem{CO_2}，也可能生成一氧化碳 \chem{CO}）。
\item \textbf{反应没走完}：可逆反应到了某个程度会“顶住”，正逆两个方向互相拉扯，剩下的反应物永远转化不完——这就是第 31 章要讲的平衡。
\item \textbf{损耗}：转移溶液时杯壁挂一层，过滤时滤纸留一点，结晶时母液里溶着一部分……每一步操作都在漏。
\item \textbf{原料不纯}：第 3 章说过，纯净物很难“纯”，称出来的 \ud{10}{g} 原料里可能只有 \ud{9.5}{g} 真正参加反应。
\end{itemize}

所以产率是化学家的KPI：产率从 60\% 提到 90\%，意味着同样的原料多出一半产品、少排一半废料。绿色化学的全部精打细算，起点都是这本账。

工业上把化学计量用到了极致。举一个有生命危险的例子：汽车安全气囊。碰撞瞬间，气囊里的叠氮化钠 \chem{NaN_3} 被电信号引爆：

\[\chem{2NaN_3(s)} \rightarrow \chem{2Na(s)} + \chem{3N_2(g)}\]

设计目标：在约 0.03 秒内生成大约 \ud{70}{L} 氮气，把气囊充满。常温下 \ud{70}{L} 氮气约 \ud{3}{mol}，按系数 $2:3$ 反推，需要 \ud{2}{mol} \chem{NaN_3}，即大约 \ud{130}{g}——多一点气囊会过压伤人，少一点起不到缓冲，药量就是这样从方程式里一格一格算出来的。（生成的金属钠活泼危险，配方里还安排了别的成分立刻把它“接住”变成无害物质——那也是按账本配的。）

\section{这本账不记什么}

化学方程式很有用，但它的边界同样清楚。学会一本账本，也要知道它\textbf{不记}什么：

\begin{itemize}[itemsep=2pt]
\item \textbf{不记能量}。方程式说氢气和氧气生成水，却一个字没提这场反应放出多少热、为什么要先点火。能量有另一本账，那是第 27、28 章的内容。
\item \textbf{不记快慢}。配平得再漂亮的方程式，也不告诉你反应在一秒内完成还是一千年。纸在空气里“应该”燃烧（热力学上完全合理），可它能安安静静放几十年——速率是第 29 章的问题。
\item \textbf{不记路径}。从反应物到生成物，中间可能经过好几步看不见的中转。方程式只记“期初余额”和“期末余额”，中间的流水被合并了。催化剂正是通过改变中间流水来加速，却不改期初期末（第 30 章）。
\item \textbf{不保证发生}。方程式是“允许发生的配比”，不是“必然发生的承诺”。能不能发生、朝哪个方向走、走多远，要靠能量和平衡来判断（第 28、31 章）。
\end{itemize}

一句话：方程式管“多少”，不管“能不能、快不快、怎么走”。把它当成万能预言书，是符号层最常见的越界。

\section{再看那支蜡烛}

回到开场的悬念。蜡去哪了？蜡的主要成分是碳氢化合物（比如二十几个碳的烷烃），燃烧时和氧气重新组合，碳原子走进二氧化碳，氢原子走进水蒸气，全部分散到你周围的空气里——一个原子也没少，只是换了搭档。把蜡烛、火焰和周围的空气看成一本完整的账：烧掉多少克蜡，就有多少克氧气和它一起变成了多少克二氧化碳加水蒸气。蜡烛“消失”，只是账本的边界画得太小。

铁钉变重，是氧气入了账；气球鼓起而总质量不变，是气体没出账。你现在手里这套记账法——原子守恒、配平、系数即配比、限制反应物、理论产率——足够给任何一场化学反应立账了。

再往远处看一步：蜡烛里的碳原子并没有“消失”，它们此刻正以二氧化碳的形式飘在空气里，明年也许被一片叶子捕获，变成葡萄糖的一部分，几年后可能出现在某个人的面包里。第 16 章说过，元素在地球上周而复始地循环，从不凭空产生或消失——每一本化学反应的账本，都只是这条巨大循环账上的一个小格子。等到第三册，你会看到生命本身就是一位精打细算的记账员：细胞里的每一次呼吸、每一片叶子的光合作用，都要按配平好的方程式一粒一粒地数原子。

\section{账本之外还有能量}

不过，细心的你可能已经憋着一个问题了：氢气和氧气明明“配平就能反应”，为什么要点火才炸？蜡烛“应该”烧完，为什么不点就不烧？断键要花钱，成键能赚钱——这笔能量收支到底怎么记？

账本的下一页，记的不是原子的数目，而是能量的进出。第 27 章，我们来算这本账。

\begin{tiaozhan}
配平其实是一道代数题。拿乙烷燃烧来说：设方程式为 $a\chem{C_2H_6} + b\chem{O_2} \rightarrow c\chem{CO_2} + d\chem{H_2O}$。列守恒方程：碳 $2a=c$，氢 $6a=2d$，氧 $2b=2c+d$。令 $a=1$，立刻得到 $c=2$，$d=3$，$b=(2c+d)/2=7/2$。出现分数没关系，全部乘以 2 即可：$\chem{2C_2H_6} + \chem{7O_2} \rightarrow \chem{4CO_2} + \chem{6H_2O}$。更一般地，含碳、氢、氧的燃料 \chem{C_xH_yO_z} 完全燃烧的通式是
\[\chem{C_xH_yO_z} + \left(x+\tfrac{y}{4}-\tfrac{z}{2}\right)\chem{O_2} \rightarrow x\chem{CO_2} + \tfrac{y}{2}\chem{H_2O}\]
不信就数一遍原子验证。这套“设未知数、列守恒、取整”的办法能配平绝大多数方程式；涉及电子转移的反应，还要再加一条“电子守恒”，第 33 章会发给你那把新尺子。跳过本节不影响主线。
\end{tiaozhan}

\section*{练一练}

\begin{sikaoti}
\item 铁在氧气中点燃生成四氧化三铁 \chem{Fe_3O_4}。先写出未配平的方程式，再数原子配平，最后验算：若 \ud{168}{g} 铁完全反应，需要多少克氧气？生成多少克 \chem{Fe_3O_4}？（相对原子质量：Fe $=56$，O $=16$）
\item 画一幅粒子图：用两种颜色的小圆代表氢原子和氧原子，画出“$2\chem{H_2} + \chem{O_2} \rightarrow \chem{2H_2O}$”反应前后的粒子排列。数一数图里每种原子的总数，并在图旁注明：圆只是表示法，原子不是硬邦邦的小球。
\item 继续做三明治：每份需要 2 片面包、1 片奶酪、1 片火腿。现有 12 片面包、5 片奶酪、7 片火腿。最多能做几份？哪种原料是“限制反应物”，各剩多少？把这个逻辑对照到第 4 题的化学情境。
\item 锌与稀盐酸反应：$\chem{Zn(s)} + \chem{2HCl(aq)} \rightarrow \chem{ZnCl_2(aq)} + \chem{H_2(g)}$。若有 \ud{6.5}{g} 锌（\ud{0.1}{mol}）和含 \ud{0.1}{mol} \chem{HCl} 的盐酸，理论能生成多少摩尔氢气？谁是限制反应物？反应结束后烧杯里还剩什么？
\item 第 4 题的实验里，某同学实际收集到的氢气只有理论值的 90\%。列出至少三个可能的原因，并分别说明它们属于“副反应、没反应完、操作损耗、原料不纯”中的哪一类。
\item 碳酸氢铵化肥受热易“失踪”：$\chem{NH_4HCO_3(s)} \xrightarrow{\text{受热}} \chem{NH_3(g)} + \chem{H_2O(g)} + \chem{CO_2(g)}$。一位农民发现一袋化肥敞口放久了“变轻了”，认为质量守恒在这里失效了。用“账本边界”的想法向他解释真相，并说明怎样设计一个实验证明质量依然守恒。
\end{sikaoti}
